% §1 Introduction — target 1 page (~5-6 paragraphs)
% Purpose: plant the network-scale gap and the 4 contributions.
% Do: lead with a concrete IoT-at-scale incident; pose the problem; state the gap; announce contributions.
% Don't: hardware tools, "autonomous exploitation", "graph-guided".

\section{Introduction}
\label{sec:intro}

% ── ¶1 hook: real-world IoT pentest motivation ────────────────────────────
% Concrete incident referencing mass IoT compromise (Mirai variants, smart-city
% camera takeovers, ICS reconnaissance campaigns). Frame as: attacker does not
% compromise *a* device — they compromise *a network* of devices.
\todo{Write ¶1: IoT attacks are network-scale; defenders pentest with single-host tooling.}

% ── ¶2 the promise, the cost ──────────────────────────────────────────────
% Manual IoT pentest is expensive and slow; LLM agents (PentestGPT~\cite{pentestgpt24},
% VulnBot~\cite{vulnbot25}, AutoPentester~\cite{autopentester25}, CAI~\cite{cai25})
% promise automation and have shown strong results on CTF and HackTheBox targets.
\todo{Write ¶2: LLM agents are progressing fast, cite 4-5 recent peers.}

% ── ¶3 the gap ────────────────────────────────────────────────────────────
% Every published LLM-pentest benchmark — AUTOPENBENCH~\cite{autopenbench24},
% HackTheBox, VulnHub, CAIBench/RCTF2~\cite{caibench25} — uses an *isolated host*
% as the unit of evaluation. Real IoT pentest is a network task: heterogeneous
% devices, multiple protocols (MQTT, Modbus, LoRaWAN, HTTP, SSH), lateral reachability
% constrained by architecture. There is no benchmark that evaluates LLM agents at
% this unit.
\todo{Write ¶3: formalize the structural gap. Cite 5-6 benchmarks. Keep it sharp.}

% ── ¶4 contribution statement ─────────────────────────────────────────────
\noindent\textbf{Contributions.} We bridge this gap with:
\begin{itemize}
  \item \textbf{\systemname} — the first reproducible benchmark for LLM pentest
  agents at \emph{network scale}: 5 architectural patterns, 7~scenarios, 8~protocols,
  87~vulnerabilities deterministically injected via Ansible, with ground truth mapped
  to OWASP~IoT~Top~10 and MITRE ATT\&CK for ICS.
  \item \textbf{\harnessname} — a network-native harness with \emph{discovery $\to$
  per-device parallel triage $\to$ per-vulnerability verification $\to$ network-level
  report}, designed to keep LLM context tractable at $/24$ scale.
  \item \textbf{Evidence-level scoring (L1/L2/L3)} — a severity-weighted evaluator
  that distinguishes detection, exploitation, and data exfiltration, beyond the
  binary flag-capture of single-host benchmarks.
  \item \textbf{Empirical study} — using a single shared LLM (\texttt{MiniMax-M2.7})
  to isolate architectural contributions, we compare \harnessname\ against
  three categories of baselines: a non-LLM scanner (Nmap NSE), single-agent
  LLM (PentestGPT), and multi-agent LLM systems with and without graph guidance
  (CAI in two scope variants, VulnBot with task-graph). We additionally report
  two ablations of our pipeline that isolate the individual contributions of
  the infrastructure-graph phase and the multi-hop intrusion phase.
\end{itemize}

% ── ¶5 findings teaser (with numbers) ─────────────────────────────────────
\todo{Write ¶5: one teaser sentence per RQ. e.g. "Best LLM reaches XX\% recall at
/24 scale; scanners reach YY\%; but only ZZ\% of LLM findings achieve L2+."}

% ── ¶6 paper organization ─────────────────────────────────────────────────
\noindent\textbf{Paper organization.}
\Cref{sec:related} surveys LLM pentest agents and benchmarks.
\Cref{sec:threat} formalizes the threat model.
\Cref{sec:iotbench} presents \systemname.
\Cref{sec:harness} describes the harness architecture.
\Cref{sec:setup} and \Cref{sec:eval} present experiments and results.
\Cref{sec:discussion} and \Cref{sec:ethical} discuss limitations and ethics.
Artifacts are released at \texttt{ANONYMIZED\_URL}.
